\documentclass[a4paper,12pt]{article}
\usepackage[utf8]{inputenc}
\usepackage[T1]{fontenc}
\usepackage{amsmath}
\usepackage{booktabs}
\usepackage{hyperref}
\usepackage{listings} % For code snippets
\usepackage{xcolor} % For code highlighting

% Configuring packages for compatibility
\usepackage{geometry}
\geometry{margin=1in}

% Adding comprehensive preamble
\usepackage{fancyhdr}
\pagestyle{fancy}
\fancyhf{}
\fancyhead[L]{ECG Arrhythmia CNN Report}
\fancyhead[R]{Wajih Humayun Hashmi, May 17, 2025}
\fancyfoot[C]{\thepage}

% Code listing setup
\lstset{
    language=Python,
    basicstyle=\ttfamily\small,
    keywordstyle=\color{blue}\bfseries,
    stringstyle=\color{red},
    commentstyle=\color{green!50!black},
    numbers=left,
    numberstyle=\tiny,
    stepnumber=1,
    numbersep=5pt,
    breaklines=true,
    frame=single,
    tabsize=4
}

% Font configuration at the end
\usepackage{times}

\begin{document}

\begin{titlepage}
    \centering
    \vspace*{1cm}
    {\Huge ECG Arrhythmia CNN Report}\par
    \vspace{0.5cm}
    {\Large Wajih}\par
    \vspace{0.5cm}
    {\large May 17, 2025}\par
    \vspace{2cm}
\end{titlepage}

\section*{Introduction}
This report details the development of a Convolutional Neural Network (CNN) for ECG arrhythmia detection using the MIT-BIH Arrhythmia Database, achieving a balanced accuracy of ~0.88 on records 100, 101, and 203.

\section{Methodology}

\subsection{Data Collection and Preprocessing}
Records 100, 101, and 203 were used, each with 40,000 samples, totaling 120,000 samples. The \texttt{ECGDataLoader} class handles loading and preprocessing:
\begin{itemize}
    \item \textbf{load\_ecg\_data}: Loads signals from CSV files and annotations from .atr files, concatenating signals and adjusting R-peak offsets.
    \item \textbf{segment\_beats}: Segments signals into 256-sample windows around R-peaks, filtering out invalid segments and assigning binary labels (0 for normal, 1 for arrhythmic).
    \item \textbf{preprocess\_data}: Normalizes segments to [0, 1] and reshapes for CNN input.
\end{itemize}

Here’s the core of \texttt{segment\_beats}:
\begin{lstlisting}
def segment_beats(self, signal, r_peaks, labels):
    half_window = self.window_size // 2
    segments = []
    segment_labels = []
    for i in range(len(r_peaks)):
        r_peak = r_peaks[i]
        start = r_peak - half_window
        end = r_peak + half_window
        if 0 <= start and end < len(signal):
            segment = signal[start:end]
            segments.append(segment)
            segment_labels.append(1 if labels[i] == 'A' else 0)
    return np.array(segments), np.array(segment_labels)
\end{lstlisting}
This method ensures segments align with labels, addressing earlier mismatches (e.g., 20 vs. 66 samples).

\subsection{Model Architecture and Training}
The \texttt{ECGCNNModel} class defines a CNN with convolutional layers, max-pooling, dropout, and dense layers, compiled with binary cross-entropy loss. The training process in \texttt{main.py} splits data into 70\% training, 15\% validation, and 15\% test sets, with stratification to maintain class balance:
\begin{lstlisting}
X_train, X_temp, y_train, y_temp = train_test_split(
    segments, segment_labels, test_size=0.3, random_state=42, stratify=stratify
)
stratify_second = y_temp if all(count >= 2 for count in counts_temp) else None
X_val, X_test, y_val, y_test = train_test_split(
    X_temp, y_temp, test_size=0.5, random_state=42, stratify=stratify_second
)
\end{lstlisting}
The dynamic stratification avoids errors when a class has fewer than 2 samples (e.g., only 1 arrhythmic beat in \texttt{y\_temp}).

\subsection{Visualization}
The \texttt{ECGPlotter} class plots ECG segments and confusion matrices, aiding in visual validation of model predictions.

\section{Results}
\begin{table}[h]
    \centering
    \begin{tabular}{lc}
        \toprule
        Metric & Value \\
        \midrule
        Accuracy & 0.97 \\
        Precision & 0.88 \\
        Recall & 1.00 \\
        F1 Score & 0.94 \\
        \bottomrule
    \end{tabular}
    \caption{Final Model Performance}
\end{table}
Class distribution improved from 264:2 to 424:51 with record 203, enabling effective learning of arrhythmic patterns.

\section{Challenges and Solutions}
\begin{itemize}
    \item \textbf{Imbalance}: Initially, only 2 arrhythmic beats were present; adding record 203 increased this to 51.
    \item \textbf{Split Mismatch}: Resolved by adjusting stratification logic to use subset labels (\texttt{y\_temp}).
    \item \textbf{Learning Minority Class}: The CNN learned arrhythmic patterns after increasing minority class samples.
\end{itemize}

\section{Conclusion and Future Work}
The model achieved a balanced accuracy of ~0.88, exceeding the ~0.85 target. Future enhancements could include class weights, adding records 207 and 208, or deeper architectures.

\end{document}